\documentclass{ceurart}

\usepackage{listings}
\usepackage{tabularx}
\newcommand{\term}[1]{\texttt{#1}}
\lstset{basicstyle=\small\ttfamily,breaklines=true,columns=fullflexible,
  keepspaces=true,showstringspaces=false,frame=single,aboveskip=6pt,belowskip=6pt}
\emergencystretch=1em
\sloppy

\begin{document}

\copyrightyear{2027}
\copyrightclause{Copyright for this paper by its authors.
  Use permitted under Creative Commons License Attribution 4.0
  International (CC BY 4.0).}
\conference{SWAT4HCLS 2027, February 1--4, 2027, Basel, Switzerland}

\title{VCF Core: A Semantic Foundation for RDF Representations of VCF}

\author[1,2]{Elias Crum}[
orcid=0009-0005-3991-754X,
email=elias.crum@ugent.be
]
\cormark[1]
\author[2]{Bart Buelens}[
orcid=0000-0001-7734-3747,
email=bart.buelens@vito.be
]
\author[2]{Gökhan Ertaylan}[
orcid=0000-0001-5602-6435,
email=gokhan.ertaylan@vito.be
]
\author[1]{Ruben Taelman}[
orcid=0000-0001-5118-256X,
email=ruben.taelman@ugent.be
]
\address[1]{IDLab, Department of Electronics and Information Systems, Ghent University -- imec, Belgium}
\address[2]{Flemish institute for Technological Research (VITO) Mol, Belgium}
\cortext[1]{Corresponding author.}


\begin{abstract}
% Context: Medical interpretation depends on connecting observations to knowledge.
Interpreting genomic variants for medical use requires connecting variant observations to clinical information and biomedical knowledge while retaining the evidence needed to assess them.
% Need: VCF carries identifiers and annotations, but lacks a native RDF graph model.
The Variant Call Format (VCF) carries much of this evidence, but does not provide a native graph model for linking its records, annotations and sample data to external resources.
% Task: Target the breadth of the format without claiming coverage beyond the evidence.
We address this need by developing a shared semantic target for RDF conversion across the breadth of VCF.
% Object: Identify the vocabulary and the contribution of this paper.
This paper presents VCF Core, its design rationale and its specification-based coverage assessment.
The vocabulary preserves file-specific interpretation context and supports expanded sample calls for single-sample analysis alongside condensed, sample-ordered vectors for large multi-sample files.
% Findings: Distinguish measured representation coverage from complete conformance.
A curated VCF 4.5 inventory, scored under a rubric that separates preserving source information from exposing it structurally, identifies dedicated queryable representations for all 104 logical-model entries; every entry carries executed evidence and 87 are additionally enforced by a validation rule.
The specification's two byte-level requirements are reported as a separate axis and verified on the bytes, because representing them in RDF would assert a property of a source the graph no longer holds.
The artifacts also provide all 122 reserved INFO and FORMAT declaration entries and version-specific registries and validation profiles for VCF 4.1--4.5.
% Conclusion: Explain why these findings matter to the audience.
Together, these contributions make representation choices and coverage gaps inspectable, establishing a common basis for comparing converters and linking observations without discarding their source context.
% Perspectives: State the investigations enabled, rather than claiming their results.
Future work will evaluate conversion fidelity, workload-dependent representation costs and reference-aware links to biomedical knowledge graphs.
\end{abstract}


\begin{keywords}
Variant Call Format\sep RDF\sep Genomic data\sep Vocabulary\sep SHACL
\end{keywords}
\maketitle

\section{Introduction and related work}

Genomic variant interpretation depends on relationships that extend beyond a variant call: its reference sequence, its possible functional consequences, and the biomedical knowledge against which it is assessed. Connecting these resources is therefore relevant to both research and medical uses of genomic data. Semantic Beacons illustrates this direction through federated queries over genomic variants and public knowledge graphs~\cite{bodrug2025semantic}; the GA4GH Variation Representation Specification (VRS) addresses the complementary need for computable variation descriptions and federated identification~\cite{wagner2021vrs}. Such integration requires not only identifying a biological alteration, but also retaining the context of the observation being interpreted.

The Variant Call Format (VCF) is a widely used interchange format for genomic variant data~\cite{danecek2011vcf}. It already accommodates database identifiers, reference information and extensible annotations. Its limitation for Linked Data is therefore not an absence of links in principle: it lacks a native RDF graph model through which those elements and their interpretation context can be addressed and queried consistently. A record's meaning depends on more than position and alleles. Header declarations define annotations, FORMAT determines the contents of sample fields, and file metadata supplies reference and processing context~\cite{vcf45}. Converting only selected variant attributes leaves downstream consumers to reconstruct these relationships or accept their loss.

Genotype data adds a second design problem. For a single-sample VCF, explicit sample calls and typed depth values support focused inspection and linking. In a large multi-sample VCF, materializing each sample/field intersection creates graph resources across the product of records, samples and FORMAT fields. Both workloads need the same source semantics, but need not expose every value at the same graph granularity. Making that choice part of a shared vocabulary allows converter and query designs to be compared against a common target.

Existing work provides substantial foundations. GFVO harmonizes GFF3, GTF, GVF and VCF concepts~\cite{baran2015gfvo}, and VCF2RDF describes an isomorphic VCF-to-RDF mapping~\cite{penha2017vcf2rdf}. HERO-Genomics models genomic entities, including VCF files and records, for integration with wider biomedical information~\cite{cazzaro2025hero}. The gvar schema concentrates on genomic variants as Linked Data and explicitly distinguishes this from representing VCF itself~\cite{gvar}. The Sequence Ontology (SO) and FALDO supply reusable descriptions of sequence features and locations~\cite{eilbeck2005so,bolleman2016faldo}. Our earlier conversion-framework paper proposed an ontology for single-sample VCF transformation~\cite{crum2025semantifying}. These approaches motivate a VCF-specific contribution with an explicit account of representation coverage, version-dependent rules and alternative genotype granularities.

We present VCF Core Vocabulary 2.0.0~\cite{vcfc200}, an independently specified semantic target developed from the VCF-RDFizer Vocabulary. Our aim is to cover the breadth of VCF while making the achieved scope testable. The contributions are: a model that preserves source context and header-defined meaning while reusing external vocabularies (Section~\ref{sec:design}); complementary expanded and condensed genotype representations (Section~\ref{sec:profiles}); a construct-by-construct VCF 4.5 coverage assessment and version-specific artifacts for VCF 4.1--4.5 (Section~\ref{sec:coverage}); and implementation and validation evidence that establishes a basis for subsequent conversion and linking studies (Section~\ref{sec:implementation}). Throughout, \term{vcfc:} denotes \url{https://w3id.org/vcf-core/vocab#}.

\section{Scope and design rationale}
\label{sec:design}

\subsection{Preserving the observation and its source context}

VCF Core represents a file's observations rather than assuming that each row identifies a globally normalized biological variant. It separates a \term{VCFFile}, its \term{VCFHeader} and header lines, individual \term{VCFRecord} resources, and a \term{VariantCall} associated with each record. The call groups QUAL, FILTER, INFO, FORMAT and genotype data; it is a modeling abstraction, not an additional VCF line type. Figure~\ref{fig:model} summarizes the shared structure and the genotype alternatives.

This separation matters when observations are linked. Two files can describe an equivalent alteration yet report different quality scores, filters, annotations or sample calls. Merging their records solely on variant identity would also merge evidence that downstream interpretation may need to distinguish. The file is therefore modeled as a \term{dcat:Distribution} and a \term{prov:Entity}; records, calls and sample resources also provide provenance attachment points~\cite{prov2013}. Recommended file-based IRI templates retain the connection between source, record, call and field. Producers must choose unambiguous record identifiers, because neither position alone nor a potentially missing VCF ID guarantees row identity.

The property \term{asSequenceAlteration} provides a separate link from a record or call to an SO sequence alteration (\term{SO:0001059}). This permits the biological description and its links to evolve while the source observation remains available for inspection. Establishing such a link requires reference and allele interpretation; the vocabulary supplies the attachment point rather than an automatic normalization procedure. Similarly, a \term{VCFSample} identifies a sample column within a file without asserting that the column is a biological individual. Pooled samples, synthetic columns and tumour/normal analyses make that distinction consequential for later links to clinical entities.

\subsection{Reusing domain models at the appropriate strength}
\label{sec:reuse}

The design favors published terms where they fit the intended meaning and introduces VCF-specific terms for the format's own structures. It distinguishes reuse in logical axioms from alignment annotations, because these have different consequences for consumers. A superclass or range contributes to the interpretation of an instance; an alignment documents a relationship between vocabulary terms without asserting that the VCF resource is identical to a resource in another model.

Logical axioms reuse five external terms. The file and provenance layers use \term{dcat:Distribution}, \term{prov:Entity} and \term{prov:wasDerivedFrom}; alteration and interval relations use \term{SO:0001059} and \term{faldo:InRangePosition} as ranges. For example, the pedigree relation corresponding to a clonal \term{Original=} declaration specializes \term{prov:wasDerivedFrom}. Such reuse connects VCF-specific structures to existing descriptions of distributions, provenance and sequence alterations while retaining their more precise role in the source format.

FALDO is also used directly in instance data~\cite{bolleman2016faldo}. Record and allele locations use \term{faldo:location}, positions or regions, and \term{faldo:reference} links to the relevant contig. Imprecise structural-variant bounds can use \term{faldo:InRangePosition}; base-modification carriers reuse FALDO's stranded-position classes. The vocabulary consequently introduces no competing position, region or strand class. This allows consumers to reuse FALDO location patterns, although cross-file comparison still depends on identifying compatible reference sequences.

The ontology additionally records 45 distinct external alignment targets: SO~(15), ChEBI~(10), VRS~(8), HERO-Genomics~(6), FALDO~(3) and GENO~(3). These occur in SKOS mapping annotations or \term{rdfs:seeAlso}, rather than as equivalent-class axioms. For example, symbolic structural-variant allele types can be related to SO classes, and base-modification aliases to the ChEBI identifiers named by the VCF specification. The distinction preserves a useful boundary: similarities between models can guide future integration without making every source record depend on an interpretation that its file has not established. These alignments identify candidate correspondences; successful instance-level joins require separate evaluation.

\begin{figure}[t]
\centering\small
\fbox{\parbox{0.92\linewidth}{\centering
\textbf{Shared VCF context}\\[1pt]
\term{VCFFile} $\rightarrow$ \term{VCFRecord} $\rightarrow$ \term{VariantCall};
\term{VCFHeader} contains header lines linked to field definitions}}
\par\smallskip
\begin{tabularx}{\linewidth}{@{}>{\centering\arraybackslash}X >{\centering\arraybackslash}X@{}}
$\downarrow$ \textbf{Expanded} & $\downarrow$ \textbf{Condensed} \\
\term{SampleCall} & \term{CohortCallMatrix} \\
$\downarrow$ \term{hasFormatValue} & $\downarrow$ \term{hasFormatValueVector} \\
\term{FormatFieldValue} & \term{FormatValueVector} \\
One resource per sample/key & One vector per key, ordered by \term{SampleSet}
\end{tabularx}
\caption{Selected conceptual structure. Both genotype branches attach to a \term{VariantCall} and retain the same file and record context. Expanded values can link to FORMAT definitions through \term{declaredBy}; condensed vectors must. The file declares the chosen representation.}
\label{fig:model}
\end{figure}

\subsection{Keeping header-defined meaning queryable}

Header declarations are part of the representation because they govern the interpretation of values. All eleven meta-information line types defined by VCF 4.5 have dedicated classes, as does the \term{\#CHROM} column line. A generic \term{HeaderAttribute} additionally carries key-value attributes of structured lines, accommodating producer-specific declarations without a new predicate for each. INFO and FORMAT definitions expose \term{fieldId}, \term{fieldNumber}, \term{fieldType} and \term{fieldDescription}, together with \term{fieldSource} and \term{fieldVersion}; field values and condensed vectors can retain the connection to their definitions through \term{declaredBy}.

A specification-level definition and its occurrence in a file are distinct resources. The reserved INFO/DP declaration, for example, describes a standardized key, whereas a \term{\#\#INFO} line records what a particular file declares. Accordingly, \term{FieldDefinition} is separate from \term{HeaderLine}, with \term{declaredOnLine} linking them. This prevents a shipped registry entry from being mistaken for a line in every source file and lets a consumer inspect the declaration actually associated with an observation.

The generic definition/value pattern also keeps identically named keys distinct when their scope or declaration differs: INFO/DP and FORMAT/DP need not describe the same measurement. It provides a consistent way to retrieve local semantics, while harmonization across producers remains an additional mapping task. Cardinality illustrates why this context is needed. \term{fieldNumber} retains the VCF string, including symbolic values whose interpretation depends on alleles or genotypes; \term{fieldArity} links to one of nine arity individuals and \term{fieldNumberInteger} represents fixed counts. Consumers can thus identify the cardinality category without first parsing the declaration.

\subsection{Separating lexical preservation from parsed interpretation}

Lexical and parsed values serve complementary purposes. \term{fieldValue} can retain a source token or comma-separated list, while optional integer, decimal and boolean properties expose suitable scalar interpretations. Numeric queries then become possible without coercing every token into an unsuitable datatype: non-finite VCF floating-point values, for example, cannot be represented as \term{xsd:decimal}. Optional raw properties retain additional source information, including INFO and FORMAT order and condensed sample blocks.

Missingness and ordering require explicit treatment because they affect interpretation. A standalone missing value can use \term{"."\textasciicircum{}\textasciicircum{}vcfc:Null}, whereas a genotype such as \term{./.} retains its internal separators and a missing vector cell retains its sample position. A missing depth must not become zero, and an unapplied FILTER must remain distinguishable from PASS. Likewise, \term{Number=.} is a cardinality declaration, not a missing observation. Ordered allele resources use \term{alleleIndex}=0 for REF and positive indices for ALT in source order, supplying the context for allele-indexed values and genotype interpretation. Expanded values may still contain lexical arrays: expansion refers to sample and field granularity, not compulsory decomposition of every token.

\section{Two complementary genotype representations}
\label{sec:profiles}

\subsection{Expanded access for single-sample analysis}

With \term{vcfc:ExpandedRepresentation}, a call links through \term{hasSampleCall} to a \term{SampleCall} for each represented sample. Each sample call identifies the sample and links through \term{hasFormatValue} to individual \term{FormatFieldValue} resources. An optional \term{forSample} relation connects it to a reusable sample-column identity.

For single-sample files, this structure makes genotype-level access a direct graph traversal. Retrieving calls with read depth at least 20 joins a FORMAT value to the definition whose \term{fieldId} is DP and filters its \term{fieldValueInteger}. Individual calls and values can also receive links and annotations. These addressable resources are useful when a downstream application needs to inspect or link a specific sample observation. The same profile remains applicable to multi-sample files when direct cell access justifies its materialization cost; no experimentally established sample-count threshold is prescribed.

\subsection{Condensed access for multi-sample cohorts}

The condensed profile factors sample-column identity into one \term{SampleSet}, whose \term{VCFSample} members have an exact \term{sampleName} and a one-based \term{sampleIndex}. A genotype-bearing call links through \term{hasCallMatrix} to a \term{CohortCallMatrix}. This matrix identifies the sample set and groups one \term{FormatValueVector} per represented FORMAT key. Each vector declares its field definition, \term{valueEncoding} and \term{encodedValues} payload.

The initial encoding, \term{VCFTextVector}, stores tab-separated raw FORMAT values in sample-index order. Every represented sample occupies one position, including missing values. Commas within a value remain part of that value, preserving arrays such as allelic depths without confusing them with sample boundaries. The call's \term{formatRaw} retains key order, and \term{sampleDataRaw} can preserve the original sample blocks. The representation therefore retains sample values rather than replacing them with cohort summaries.

This design reduces the number of separately materialized graph resources while retaining explicit file, record, field and sample-order context. It also makes the access cost visible: an RDF literal containing a vector does not expose its cells through ordinary basic graph patterns. Per-sample comparisons require decoding through application logic or query functions. The two profiles thus offer a deliberate choice between direct access to individual values and less graph structure around the same values. Future selective expansion could combine these access patterns, provided that redundant representations are kept consistent.

\subsection{Illustration and structural consequences}

Consider a synthetic biallelic record at position 100 on contig 1, with REF=A, ALT=G and FORMAT=GT:DP. Its three sample blocks are \term{0/1:42}, \term{0/0:18} and \term{./.:.}. The condensed representation has the following payloads, where \term{\textbackslash t} denotes a tab:

\begin{lstlisting}
GT: "0/1\t0/0\t./."
DP: "42\t18\t."
\end{lstlisting}

Selecting position two in each vector recovers SAMPLE2's genotype \term{0/0} and depth 18; the third depth remains missing. An equal-content expanded graph has three sample-call resources and six FORMAT-value resources; a single-sample projection has one sample call and two values. The expanded depth pattern selects SAMPLE1, and the condensed representation gives the same selection after decoding and excluding the missing depth. The accompanying executable check recovers all six cells and verifies this answer equivalence. It demonstrates the encoding on this example, not universal conversion fidelity.

The structural distinction can be stated independently of a storage engine. Assuming $R$ genotype-bearing records each representing the same $S$ samples and $F$ FORMAT keys, expanded materialization creates $RS$ sample calls and $RSF$ field-value resources. Condensed materialization creates $S$ sample resources, one sample set, $R$ matrices and $RF$ vectors. Excluding shared file, header and record structures, the counts are
\[
N_{\mathrm{expanded}}=RS+RSF,\qquad
N_{\mathrm{condensed}}=S+1+R+RF.
\]
The optional reusable sample layer can also be added to the expanded profile. These counts explain why condensation is relevant to large cohorts, but the payload still contains $RSF$ cells. They establish a difference in graph materialization, not a measured reduction in serialized bytes or query time. Those outcomes depend on encoding, indexing and implementation.

\section{Specification coverage and version-specific validation}
\label{sec:coverage}

\subsection{Assessment method and scope}

To make the intended breadth of VCF coverage assessable, we constructed an inventory by reviewing the VCF 4.5 specification section by section~\cite{vcf45}. The inventory contains 104 entries across seven areas (Table~\ref{tab:coverage}). Each entry is scored on three independent axes rather than by free judgement: \emph{preserved} asks whether the source information survives in the graph at all, \emph{structured} asks whether the construct is reachable by a graph pattern without parsing a literal, and \emph{enforced} asks whether a validation rule tests a stated requirement. An entry counts as fully represented only when it is both preserved and structured, so a retained raw string is never credited as structure; the report recomputes each verdict from its axes and fails if the two disagree. Each entry also records its strongest evidence---an executed query with a fixed expected answer, a negative probe showing the governing rule rejects a violating graph, a validation rule, or a materialized fixture. It is a published, curated operationalization of coverage; it is not an independently established complete enumeration of every normative requirement.

Each entry identifies the specification section, the relevant vocabulary terms, supporting artifact references and a coverage verdict. \emph{Full} means that a dedicated queryable term or term set exists for the construct. \emph{Partial} means that representation relies on a documented convention or retained raw values rather than a dedicated structured expression. \emph{Absent} means that the inventory identifies no carrier. These verdicts assess representational support, not whether every constraint on that construct has been implemented or every possible VCF input tested. Reserved INFO and FORMAT declarations are counted separately: the ability to express a field definition and the inclusion of particular standardized definitions are different contributions.

\begin{table}[t]
\caption{Coverage of the 104 logical-model entries in the curated VCF 4.5 inventory. Byte-level requirements are assessed on a separate axis (Section~\ref{sec:serialization}); reserved key declarations are assessed separately again.}
\label{tab:coverage}
\centering\small
\begin{tabularx}{\linewidth}{@{}Xrrrr@{}}
\hline
Specification area & Constructs & Full & Partial & Absent \\
\hline
File and header structure & 20 & 20 & -- & -- \\
Field declaration attributes & 8 & 8 & -- & -- \\
Type and arity value sets & 14 & 14 & -- & -- \\
Lexical and encoding rules & 7 & 7 & -- & -- \\
Fixed fields and allele structure & 18 & 18 & -- & -- \\
FORMAT and genotype & 21 & 21 & -- & -- \\
Structural variation, repeats, gVCF & 16 & 16 & -- & -- \\
\hline
\textbf{Total} & \textbf{104} & \textbf{104} & \textbf{--} & \textbf{--} \\
\hline
\end{tabularx}
\end{table}

\subsection{Coverage findings and remaining gaps}

The assessment identifies 104 fully represented constructs. An earlier round of this assessment scored four of them as partial---the angle-bracketed assembly-contig form of CHROM, telomeric POS conventions (0 and $N+1$), REF padding-base conventions, and scalar versus itemwise missingness---and resolving them required two different kinds of work, which the rubric now keeps apart. Telomeric positions and missingness were already structurally representable and had been undercredited: numeric POS with a linked contig length reaches both boundaries, and a whole-field \term{Null} is distinguishable from per-item \term{Null} values on indexed field items. They needed evidence, not terms, and now carry executed queries. The assembly-contig form and the padding convention were genuine gaps and did need terms: a bracketed CHROM now resolves to an \term{AssemblyContig} carrying its parsed identifier and its sequence source, deliberately not to a fabricated \term{ContigHeaderLine}, and a \term{PaddingInterpretation} records which base serves as padding together with the specification rule that determined it. Reviewing the rubric also demoted one previously full entry before repair: dropped trailing FORMAT fields had been credited from a raw field alone, which the rubric no longer allows, so the declared FORMAT key list became a resource sequence and an omitted key is now distinguishable from an undeclared one.

The evidence distribution matters as much as the total. Every entry carries executed evidence: 7 rest on an executed query, 23 on a negative regression probe, 62 on a validation rule and 12 on a materialized fixture. 14 constructs are exercised by one of 15 queries with fixed expected answers, and 87 of 104 are enforced as well as represented. No entry rests on the existence of a term alone.

\subsection{Source-byte requirements as a separate axis}
\label{sec:serialization}

VCF 4.5 also requires the source to be UTF-8 without a byte order mark and to separate lines with LF or CR+LF. These constrain the octet stream rather than the logical model, and an RDF property for them would assert a property of a file the graph no longer holds without showing the file was read. Consistent with excluding the BCF byte layout for the same reason, the vocabulary mints no term for them; instead a byte-level checker reads every fixture, and 16 of 16 satisfy the requirements, with both permitted terminators exercised and paired LF and CR+LF fixtures shown to materialize to isomorphic graphs. Moving these two entries out of the construct denominator is the only change we have made to it, and the inventory records the change with its rationale so it cannot be mistaken for the entries quietly disappearing. The resulting claim is narrower than representation and is reported as such.

The two absent properties concern UTF-8/byte-order-mark state and LF versus CRLF line endings. They are retained in the reported denominator even though byte-for-byte reconstruction is outside the vocabulary's intended scope. BCF~2.2 binary layout is also excluded: the vocabulary targets VCF content rather than a binary encoding. These boundaries distinguish a representation intended for semantic access from a complete archival description of the source byte stream.

In addition to the construct assessment, the VCF 4.5 registry includes all 122 reserved key declarations, comprising 51 INFO and 71 FORMAT definitions, with their specified Number, Type and Description. Three entries represent the parameterized M, DPM and ADM base-modification key families, rather than individual identifiers. The registry is generated from the specification source and records the input SHA-256. Across the core, alleles, genotypes, structural-variation and reserved-key modules, the census reports 466 declared resources: 78 classes, 75 object properties, 111 datatype properties, 189 named individuals, 7 annotation properties, 5 datatypes and the ontology resource. These artifact counts describe the implementation's extent; the construct assessment supplies the more relevant evidence of coverage.

\subsection{Reproducibility and interpretation of the assessment}

The command \term{npm run coverage:report} recomputes counts from the ontology modules, SHACL profiles, versioned registries and curated inventory. It writes \term{coverage/curated/generated/report.json} and fails if the inventory cites an undeclared vocabulary term. A companion check compares selected numerical statements in the manuscript with that report. This makes the calculation and its term references inspectable and helps detect drift between artifacts and manuscript.

Automation does not establish that an editorial coverage verdict is correct, that its cited evidence is sufficient, or that the inventory is exhaustive. Those claims require review of the published entries. The assessment is useful precisely because a reader can inspect, challenge or extend an individual entry and reproduce the resulting totals. It supplies an explicit baseline for subsequent VCF-to-RDF work.

\subsection{One shared model with version-specific constructs and rules}

VCF releases share file, record and genotype structures while differing in available constructs and their permitted encodings. VCF Core represents the shared structures once and includes terms needed for later features; version-specific registries and SHACL overlays determine which declarations and constraints apply to a file. This avoids duplicating the common model while retaining distinctions that affect interpretation. It does not assume that all versions have identical expressive scope.

Five overlays, \term{vcf-4.1.shacl.ttl} through \term{vcf-4.5.shacl.ttl}, condition their constraints on the owning file's \term{fileFormat}. They can therefore be loaded together while applying version-specific rules. Optional \term{VCF41File}\ldots\term{VCF45File} classes also check agreement between the intended profile and the declared version. For example, CIPOS has two values in VCF 4.1--4.3 and two per alternate allele in VCF 4.4--4.5. A historical two-value encoding must consequently be assessed in its version context. The local-allele Number codes LA, LR and LG and the base-modification code M are restricted to VCF 4.5; their availability is not inferred merely from the presence of a vocabulary term.

The registries contain 40 explicit reserved declarations for each of VCF 4.1 and 4.2, 67 for 4.3, 78 for 4.4 and 122 for 4.5. Each snapshot records its official source URL and source SHA-256. For 4.1 and 4.2, prose-only INFO fields whose precise format is left to producers are not assigned invented Number or Type constraints. VCF 4.0 has no published overlay and is reported as unsupported, rather than assessed under a later version's rules.

These overlays complement three shared validation layers: a portable structural profile of 53 node shapes and 115 property shapes, a SPARQL profile of 16 cross-resource ordering and uniqueness constraints, and a consistency profile of 24 rules relating raw tokens to parsed values~\cite{shacl2017}. Separating the layers accommodates validators with different capabilities and makes the checks applied to a graph explicit. Future VCF changes can extend the appropriate registry or constraints and, when necessary, the model itself.

\section{Implementation evidence and implications}
\label{sec:implementation}

\subsection{Available artifacts and converter use}

VCF Core 2.0.0 provides OWL/Turtle definitions, SHACL shapes, example graphs, illustrative mapping patterns and generated documentation~\cite{vcfc200}.

The model also has an implementation lineage in VCF-RDFizer, an open-source VCF-to-RDF converter distributed under the MIT license~\cite{vcfrConverter2026}. In the cited source snapshot, default RML mappings construct file, header, record and call resources, while streaming emitters produce either per-sample calls and FORMAT values or cohort matrices and FORMAT vectors. The \term{--sample-representation} option selects the genotype structure. This snapshot uses the predecessor VCF-RDFizer namespace; it demonstrates implementation of the shared design and both genotype strategies, rather than verified conformance to every construct of VCF Core 2.0.0. The converter's documented execution paths and regression tests provide concrete implementation evidence, while migration to the renamed vocabulary and assessment against its current profiles remain separately checkable tasks.

\subsection{What validation establishes}

The recorded full validation run accepts 21 fixtures with zero violations and zero warnings, covering maintained complete examples, both paper graphs and per-version examples. A separate negative control and 35 targeted regression tests exercise invalid inputs and valid edge cases, including historical and current CIPOS rules. Decoded-value checks run in a complementary Python layer where portable SHACL is insufficient; these include genotype-likelihood cardinalities at arbitrary ploidy and checks over vector contents. The repository reports these layers separately so that the source of each check is visible.

Three forms of evidence must be distinguished. The coverage inventory assesses whether constructs have representation carriers. Constraint validation assesses graphs against the implemented rules and test cases. Paired representation checks, such as the six-cell example in Section~\ref{sec:profiles}, assess recovery of particular source values. None substitutes for the others. In particular, optional raw properties allow two valid graphs to retain different amounts of source text, and passing shapes does not establish faithful conversion of an arbitrary input file. The present corpus supports the tested cases without constituting exhaustive VCF conformance or a demonstration of downstream interoperability.

\subsection{Implications for conversion and linked-data research}

An implementer can identify the vocabulary construct corresponding to a VCF feature, inspect the applicable version rules, and choose how sample values should be exposed. A researcher can compare converters against the same expected content rather than confounding conversion behavior with undocumented model differences. The published gaps also make it possible to distinguish a vocabulary limitation from an implementation omission.

For conversion research, the next step is an equal-content fidelity suite covering phased and partially missing genotypes, variable ploidy, multiallelic and structural variants, and heterogeneous header declarations. Representation experiments can then vary records, samples, FORMAT fields and missingness while measuring triples, serialized and compressed bytes, peak memory, conversion time and loading time against VCF/BCF baselines. Query experiments should verify identical answers before comparing direct expanded access, vector decoding and selective expansion. The converter's optional graph-compression stages should be measured separately from the choice of genotype granularity.

For linked-data research, source-preserving observations provide a starting point for reference-aware links to normalized variation and external annotations. A study could connect observations to biomedical knowledge graphs while retaining the originating record, field declaration and sample context, then evaluate join accuracy and the ability to trace each answer back to its evidence. HERO-Genomics and Semantic Beacons illustrate integration settings for such investigations~\cite{cazzaro2025hero,bodrug2025semantic}. VCF Core supplies part of the representational foundation; reference reconciliation, annotation mappings and evaluation of clinical usefulness remain work to be undertaken.

\section{Conclusion}

VCF Core makes a broad VCF-to-RDF representation target explicit, inspectable and reusable. Its construct inventory establishes the achieved scope---all 104 logical-model entries fully represented, each with executed evidence, and the specification's byte-level requirements verified on a separate axis---while version-specific artifacts preserve distinctions across VCF 4.1--4.5. Expanded and condensed genotype structures address different access needs without changing the source context that gives the values meaning.

The resulting foundation enables a more precise next question than whether VCF can be expressed in RDF: which conversion and representation choices preserve the required evidence at an acceptable cost for a given workload? Answering that question against a common model can support subsequent links between genomic observations and biomedical knowledge without obscuring their provenance. The available artifacts and converter implementation lineage make these investigations concrete; conversion fidelity, performance and downstream linking are the next subjects of evaluation.

\begin{acknowledgments}
Project funding provided from the Research Foundation -- Flanders (FWO) (SB Fellowship 1S27825N).
Ruben Taelman is a postdoctoral fellow of the Research Foundation -- Flanders (FWO) (1202124N).
% KNoWS-IDLab project attribution template. Add the project name and grant
% agreement number if this work is booked to a named KNoWS-IDLab project rather
% than the first author's PhD; see ACKNOWLEDGEMENTS.md in the repository.
%   This work was carried out at KNoWS, IDLab, Ghent University -- imec, within
%   <KNoWS-IDLab project name> <grant agreement number>.
\end{acknowledgments}

\section*{Declaration on Generative AI}
OpenAI Codex and Claude assisted with literature and reference checks, outlining, drafting, the schematic and synthetic example, the vocabulary rename, the coverage inventory and measurement scripts, and LaTeX verification. The human authors must substantively review and revise this preparatory draft before submission, including its curated coverage verdicts.

% Finalize this declaration after author revision; do not claim completed human
% review prematurely. Check the current CEUR-WS Generative AI policy before submission.

% Longer working draft retained at the authors' request (8 September 2026).
% Condense the main matter to the five-page short-paper limit before submission.
% References are excluded from the five-page main-matter limit.
\clearpage
\bibliography{sources}
\end{document}
